\section{Vectors} \label{pure_vectors}
    \subsection{The Vector Product}
        For two 3D vectors, $\textbf{a}$ and $\textbf{b}$, the vector product (also called the cross product) is written $\textbf{a}\times\textbf{b}$.
        The vector product is defined as:
        \begin{equation} \label{eq:vector_product}
            \textbf{a}\times\textbf{b}=|\textbf{a}||\textbf{b}|\sin \theta\,\hat{\textbf{n}}
        \end{equation}
        Where $\theta$ is the angle between the two angles (this can be found using the dot product), and $\hat{\textbf{n}}$ is the unit vector perpendicular to $\textbf{a}$ and $\textbf{b}$.
        However, since $\textbf{a}$ and $\textbf{b}$ are both 3D, there are two opposite unit vectors that are both perpendicular to both $\textbf{a}$ and $\textbf{b}$ ($\hat{\textbf{n}}$ and $-\hat{\textbf{n}}$).

        The `right-hand rule' can be used to find the direction of $\hat{\textbf{n}}$.
        Imagining your index finger as $\textbf{a}$, your middle finger as $\textbf{b}$, the thumb of you right hand will point in the direction of $\hat{\textbf{n}}$.
        This means that swapping $\textbf{a}$ and $\textbf{b}$ will flip the direction of $\hat{\textbf{n}}$.
        Therefore, $\textbf{a}\times\textbf{b}=-(\textbf{b}\times\textbf{a})$.
        This property is called anticommutatitivity (opposed to commutatitivity, such as for addition, where $a+b=b+a$).

        Another way of thinking about flipping the direction of the vector product is by changing the direction of the angle measurement.
        $\theta$ is measured as the smallest angle between $\textbf{a}$ and $\textbf{b}$, meaning that $\sin\theta>0$.
        However, flipping $\textbf{a}$ and $\textbf{b}$ would be the same as measuring the angle in the other direction, which would give an angle $2\pi-\theta$ (in \glspl{radian}).
        Using the identity:
        \begin{equation}
            \sin\theta=-\sin(2\pi-\theta)
        \end{equation}
        Measuring the angle from the opposite direction would give a negative answer, which is the same as it would when negating $\hat{\textbf{n}}$.

        \subsubsection{Vector Product Indentities}
        From definition \ref{eq:vector_product}, the following identities can be found for the product of the basis vectors $\textbf{i}$, $\textbf{j}$, and $\textbf{k}$:
        \begin{equation} \label{eq:vector_product_of_ii_jj_kk}
            \textbf{i}\times\textbf{i}=\textbf{j}\times\textbf{j}=\textbf{k}\times\textbf{k}=0
        \end{equation}
        \begin{equation} \label{eq:vector_product_of_ij}
            \textbf{i}\times\textbf{j}=-(\textbf{j}\times\textbf{i})=\textbf{k}
        \end{equation}
        \begin{equation} \label{eq:vector_product_of_jk}
            \textbf{j}\times\textbf{k}=-(\textbf{k}\times\textbf{j})=\textbf{i}
        \end{equation}
        \begin{equation} \label{eq:vector_product_of_ki}
            \textbf{k}\times\textbf{i}=-(\textbf{i}\times\textbf{k})=\textbf{j}
        \end{equation}

        \begin{theorem}
            Using any two non-zero vectors, $\normalfont\textbf{a}$ and $\normalfont\textbf{b}$, $\normalfont\textbf{a}\times\textbf{b}=0$ if and only if they are parallel.
        \end{theorem}
        \begin{proof}
            Given definition \ref{eq:vector_product}, $|\textbf{a}\times\textbf{b}|$ can be found as the product of the magnitudes of each of of the parts of the definition:
            \begin{equation}
                |\textbf{a}\times\textbf{b}|=|\textbf{a}||\textbf{b}||\sin\theta||\hat{\textbf{n}}|
            \end{equation}
            Given that $0\le\theta\le\pi$, $\sin\theta\ge0$, so $|\sin\theta|=\sin\theta$.
            We also know that $|\hat{\textbf{n}}|$ is a unit vector by definition, so $|\hat{\textbf{n}}|=1$, so the formula for the magnitude of $|\textbf{a}\times\textbf{b}|$ simplifies to:
            \begin{equation}
                |\textbf{a}\times\textbf{b}|=|\textbf{a}||\textbf{b}|\sin\theta
            \end{equation}
            Thus, (since $\textbf{a}$ and $\textbf{b}$ are non-zero), $|\textbf{a}\times\textbf{b}|=0$ if and only if $\sin\theta=0$, which would mean that $\theta=0$ or $\theta=\pi$.
            If $\theta=0$ or $\theta=\pi$, then $\textbf{a}$ and $\textbf{b}$ must be parallel, since $\theta=0$ means that $\sin\theta=0$.
        \end{proof}

        \begin{theorem}
            Using any two vectors, $\normalfont\textbf{a}$ and $\normalfont\textbf{b}$, $\normalfont\textbf{a}\times\textbf{b}=|\normalfont\textbf{a}||\textbf{b}|$ if and only if they are perpendicular.
        \end{theorem}
        \begin{proof}
            Given definition \ref{eq:vector_product}, $|\textbf{a}\times\textbf{b}|$ can be found as the product of the magnitudes of each of of the parts of the definition:
            \begin{equation}
                |\textbf{a}\times\textbf{b}|=|\textbf{a}||\textbf{b}||\sin\theta||\hat{\textbf{n}}|
            \end{equation}
            Given that $0\le\theta\le\pi$, $\sin\theta\ge0$, so $|\sin\theta|=\sin\theta$.
            We also know that $|\hat{\textbf{n}}|$ is a unit vector by definition, so $|\hat{\textbf{n}}|=1$, so the formula for the magnitude of $|\textbf{a}\times\textbf{b}|$ simplifies to:
            \begin{equation}
                |\textbf{a}\times\textbf{b}|=|\textbf{a}||\textbf{b}|\sin\theta
            \end{equation}
            Thus, $|\textbf{a}\times\textbf{b}|=|\textbf{a}||\textbf{b}|$ if and only if $\sin\theta=1$, which would mean that $\theta=\frac{\pi}{2}$, which means that they are perpendicular, since $\frac{\pi}{2}$ is a right-angle.
        \end{proof}

        \begin{theorem}
            Where the vectors $\normalfont\textbf{a}$ and $\normalfont\textbf{b}$ are defined as follows:
            \begin{equation*}
                \normalfont{\textbf{a}}=a_1\normalfont{\textbf{i}}+a_2\normalfont{\textbf{j}}+a_3\normalfont{\textbf{k}}=
                \begin{pmatrix}
                    a_1 \\ a_2 \\ a_3
                \end{pmatrix}
            \end{equation*}
            \begin{equation*}
                \normalfont{\textbf{b}}=b_1\normalfont{\textbf{i}}+b_2\normalfont{\textbf{j}}+b_3\normalfont{\textbf{k}}
                \begin{pmatrix}
                    b_1 \\ b_2 \\ b_3
                \end{pmatrix}
            \end{equation*}
            Their vector product, $\normalfont{\textbf{a}}\times\normalfont{\textbf{b}}$, is:
            \begin{equation} \label{vector_product_ijk_form}
                \normalfont{\textbf{a}}\times\normalfont{\textbf{b}}=(a_2b_3-a_3b_2)\normalfont{\textbf{i}}+(a_3b_1-a_1b_3)\normalfont{\textbf{j}}+(a_1b_2-a_2b_1)\normalfont{\textbf{k}}
            \end{equation}
            Or in column form:
            \begin{equation}
                \normalfont{\textbf{a}}\times\normalfont{\textbf{b}}=
                \begin{pmatrix}
                    a_1 \\ a_2 \\ a_3
                \end{pmatrix}
                \times
                \begin{pmatrix}
                    b_1 \\ b_2 \\ b_3
                \end{pmatrix}
                =
                \begin{pmatrix}
                    a_2b_3-a_3b_2 \\
                    a_3b_1-a_1b_3 \\
                    a_1b_2-a_2b_1
                \end{pmatrix}
            \end{equation}
        \end{theorem}
        \begin{proof}
            We'll begin with expanding $\textbf{a}\times\textbf{b}$ using their $\textbf{i}$, $\textbf{j}$, $\textbf{k}$ forms:
            \begin{equation}
                \textbf{a}\times\textbf{b}=(a_1\textbf{i}+a_2\textbf{j}+a_3\textbf{k})\times(b_1\textbf{i}+b_2\textbf{j}+b_3\textbf{k})
            \end{equation}
            From there, the two trinomials can be expanded normally:
            \begin{equation}
            \begin{aligned}
                \textbf{a}\times\textbf{b}
                &=a_1b_1(\textbf{i}\times\textbf{i})+a_1b_2(\textbf{i}\times\textbf{j})+a_1b_3(\textbf{i}\times\textbf{k}) \\
                &+a_2b_1(\textbf{j}\times\textbf{i})+a_2b_2(\textbf{j}\times\textbf{j})+a_2b_3(\textbf{j}\times\textbf{k}) \\
                &+a_3b_1(\textbf{k}\times\textbf{i})+a_3b_2(\textbf{k}\times\textbf{j})+a_3b_3(\textbf{k}\times\textbf{k})
            \end{aligned}
            \end{equation}
            From here, we can use the identities \ref{eq:vector_product_of_ii_jj_kk}, \ref{eq:vector_product_of_ij}, \ref{eq:vector_product_of_jk}, \ref{eq:vector_product_of_ki} to simplify each term of the expanded trinomial:
            \begin{equation}
            \begin{aligned}
                \textbf{a}\times\textbf{b}
                &=a_1b_2(\textbf{k})+a_1b_3(-\textbf{j}) \\
                &+a_2b_1(-\textbf{k})+a_2b_3(\textbf{i}) \\
                &+a_3b_1(\textbf{-i})+a_3b_2(\textbf{j}) \\
                &=(a_2b_3-a_3b_1)\textbf{i}+(a_3b_2-a_1b_3)\textbf{j}+(a_1b_2-a_2b_1)\textbf{k}
            \end{aligned}
            \end{equation}
            Which is the form in equation \ref{vector_product_ijk_form}.
        \end{proof}

        This can also be written as the determinant of a $3\times3$ matrix:
        \begin{equation}
            \textbf{a}\times\textbf{b}=\left|\begin{matrix}
                \textbf{i} & \textbf{j} & \textbf{k} \\
                a_1 & a_2 & a_3 \\
                b_1 & b_2 & b_3
            \end{matrix}\right|=\textbf{i}
            \left|\begin{matrix}
                a_2 & a_3 \\
                b_2 & b_3
            \end{matrix}\right|+\textbf{j}
            \left|\begin{matrix}
                a_1 & a_3 \\
                b_1 & b_3
            \end{matrix}\right|+\textbf{k}
            \left|\begin{matrix}
                a_1 & a_2 \\
                b_1 & b_2
            \end{matrix}\right|
        \end{equation}

    \subsection{Finding Areas of Triangles and Parallelograms with Vectors}
        The area of a triangle with legs $a$ and $b$ and an angle $\theta$ between them is:
        \begin{equation}
            A=\frac{1}{2}ab\sin\theta
        \end{equation}
        This closely resembles the formula for the vector product, where $a=|\textbf{a}|$ and $b=|\textbf{b}|$.
        The using those definitions of $a$ and $b$, the formula for the area of a triangle is:
        \begin{equation}
        \begin{aligned}
            A&=\frac{1}{2}|\textbf{a}||\textbf{b}|\sin\theta \\
            &=\frac{1}{2}|\textbf{a}\times\textbf{b}|
        \end{aligned}
        \end{equation}
        Thus, this formula can be used to derive the area of a triangle from its side, where its sides are given by vectors.
        This formula can also be extended to parallelograms; the area of a parallelogram with sides $a$ and $b$ is twice the area of a triangle with sides $a$ and $b$, so its area would be given as:
        \begin{equation}
            A=|\textbf{a}\times\textbf{b}|
        \end{equation}

        \newpage

\section{Conic Sections I} \label{}

    \newpage

\section{Conic Sections II} \label{}

    \newpage

\section{Inequalities} \label{}

    \newpage

\section{The \texorpdfstring{$t$}.-Formulae} \label{}

    \newpage

\section{Taylor Series} \label{}

    \newpage

\section{Methods in Calculus} \label{}

    \newpage

\section{Numerical Methods} \label{}

    \newpage

\section{Reducible Differential Equations} \label{}

    \newpage